% !TEX root = ../main.tex
\chapter{Python--C++ Binding via pybind11}
\label{ch:python_cpp_binding}

\Lucid\ is a Python-fronted, C++-engined framework. The interface
between the two languages is the topmost layer of the engine DAG
(\chref{ch:layered_dag}), implemented through pybind11
\cite{pybind11}. This chapter describes how that binding works in
\Lucid\ specifically: the conventions, the contracts, the
exception-translation policy, the reference-counting interplay
between Python and C++ \code{shared\_ptr}, and the build-system
integration that produces the \code{lucid.\_C} extension module.

The binding layer is not large in absolute terms --- about 3500
lines of C++ across the files in \code{lucid/\_C/bindings/} --- but
it is the only place where pybind11 types appear. The discipline
the chapter describes is what keeps the engine's lower layers
pybind11-free, which in turn keeps the option open to expose the
engine through a C API to a non-Python frontend in a future
release.

\section{Overview}
\label{sec:python_cpp_binding:overview}

pybind11 is a header-only C++ library that bridges between C++
types and Python objects. The bridge has two principal operations:
\emph{type binding}, which exposes a C++ class as a Python class,
and \emph{function binding}, which exposes a C++ function as a
Python callable. Both operations are declarative: the binding code
specifies what to expose, and pybind11 generates the
boilerplate that handles argument conversion, return-value
conversion, reference counting, and exception translation.

For \Lucid\ the binding code is:

\begin{itemize}
  \item \code{lucid/\_C/bindings/bind.cpp}: the module entry point,
    \code{PYBIND11\_MODULE(engine, m)}.
  \item \code{lucid/\_C/bindings/bind\_tensor.cpp}: the
    \code{Tensor} class binding (about 400 lines).
  \item Per-subsystem binding files: \code{bind\_bfunc.cpp},
    \code{bind\_ufunc.cpp}, \code{bind\_gfunc.cpp},
    \code{bind\_composite.cpp}, \code{bind\_autograd.cpp},
    \code{bind\_nn.cpp}, \code{bind\_optim.cpp},
    \code{bind\_linalg.cpp}, \code{bind\_fft.cpp},
    \code{bind\_einops.cpp}, \code{bind\_complex.cpp},
    \code{bind\_amp.cpp}, \code{bind\_profiler.cpp},
    \code{bind\_random.cpp}, \code{bind\_compile.cpp}.
  \item \code{bind\_errors.cpp}: the exception translator.
  \item \code{bind\_op\_registry.cpp}: introspection binding for
    the op registry.
  \item \code{BindingGen.h}: helper templates for code generation.
\end{itemize}

The output of building these files is a single CPython extension
module, \code{lucid/\_C/engine.cpython-XYZ.so}, where XYZ encodes
the Python version. The module is loaded under the alias
\code{\_C\_engine} per Hard Rule H2; see
\secref{sec:python_cpp_binding:h2_alias}.

\section{Module Initialisation}
\label{sec:python_cpp_binding:module_init}

The module entry point is declared with the
\code{PYBIND11\_MODULE} macro:

\begin{lstlisting}[style=lucidcpp,language=C++]
#include <pybind11/pybind11.h>
namespace py = pybind11;

PYBIND11_MODULE(engine, m) {
  m.doc() = "Lucid C++ engine bindings";

  bind_errors(m);     // must come first --- registers exception translator
  bind_tensor(m);
  bind_ufunc(m);
  bind_bfunc(m);
  bind_gfunc(m);
  bind_composite(m);
  bind_autograd(m);
  bind_nn(m);
  bind_optim(m);
  bind_linalg(m);
  bind_fft(m);
  bind_einops(m);
  bind_complex(m);
  bind_amp(m);
  bind_profiler(m);
  bind_random(m);
  bind_compile(m);
  bind_op_registry(m);
}
\end{lstlisting}

Each \code{bind\_<subsystem>(m)} function adds its part of the
surface to the module object. The ordering matters in one place:
\code{bind\_errors} must run first so that the exception translator
is installed before any other binding might throw.

\subsection{The module-name convention}
\label{ssec:python_cpp_binding:name_convention}

The module is named \code{engine}, which means it is imported in
Python as \code{from lucid.\_C import engine}. The alias convention
of Hard Rule H2 requires that the import be aliased to
\code{\_C\_engine}:

\begin{lstlisting}[style=lucid,language=LucidPython]
from lucid._C import engine as _C_engine
\end{lstlisting}

The alias prefix \code{\_C\_} flags that the symbol is the
engine-level binding and makes greppability trivial. Section
\ref{sec:python_cpp_binding:h2_alias} treats the alias rule in
detail.

\section{Type Bindings: The Tensor Class}
\label{sec:python_cpp_binding:tensor_binding}

The \code{Tensor} class is the central exposed type. Its binding,
abbreviated:

\begin{lstlisting}[style=lucidcpp,language=C++]
void bind_tensor(py::module_& m) {
  py::class_<TensorImpl, std::shared_ptr<TensorImpl>>(m, "Tensor")
    .def(py::init<>())
    .def_property_readonly("shape",
        [](const TensorImpl& t) { return t.shape; })
    .def_property_readonly("dtype",
        [](const TensorImpl& t) { return t.dtype; })
    .def_property_readonly("device",
        [](const TensorImpl& t) { return t.device; })
    .def_property("requires_grad",
        &TensorImpl::requires_grad,
        &TensorImpl::set_requires_grad)
    .def("backward", &TensorImpl::backward,
        py::arg("gradient") = py::none(),
        py::arg("retain_graph") = false)
    .def("detach", &TensorImpl::detach)
    .def("clone",  &TensorImpl::clone)
    // ... and so on, one binding per method
    .def("__add__", [](const TensorImpl& a, const TensorImpl& b) {
        return Dispatcher::instance().binary(OpKind::Add, a, b);
    })
    // ... and the rest of the operator-overload dunders
    ;
}
\end{lstlisting}

The crucial detail is the second template argument to
\code{py::class\_}: \code{std::shared\_ptr<TensorImpl>}. This tells
pybind11 that \code{TensorImpl} instances are managed through a
\code{shared\_ptr}, which is how the engine internally tracks
tensor lifetimes. The Python-side reference counting then
participates in the same \code{shared\_ptr} count, so a tensor
stays alive as long as either the Python object or any internal
\code{shared\_ptr} refers to it.

\subsection{Property versus method}
\label{ssec:python_cpp_binding:property_method}

The \code{def\_property\_readonly} versus \code{def} choice
matches the Python convention. Attributes like \code{shape},
\code{dtype}, \code{device} are exposed as properties (no
parentheses needed at the call site), while operations like
\code{backward}, \code{detach}, \code{clone} are exposed as
methods. This is the \refframework-compatible surface that the
user expects.

\subsection{Method-like dunders}
\label{ssec:python_cpp_binding:dunders}

Operator overloads (\code{\_\_add\_\_}, \code{\_\_mul\_\_},
\code{\_\_matmul\_\_}, etc.) are bound as lambdas that call into
the dispatcher. Each lambda is one line; there are about 30 of
them. The pattern is uniform enough that we could code-generate it
from a list, but the gain is small and explicit code is easier to
diff.

\section{Function Bindings}
\label{sec:python_cpp_binding:function_binding}

Free functions are bound similarly to methods. Per-op bindings live
in \code{bind\_ufunc.cpp}, \code{bind\_bfunc.cpp},
\code{bind\_gfunc.cpp}, and \code{bind\_composite.cpp}, organised
by op category.

A typical unary binding:

\begin{lstlisting}[style=lucidcpp,language=C++]
void bind_ufunc(py::module_& m) {
  m.def("relu",    &ops::relu,    py::arg("x"));
  m.def("sigmoid", &ops::sigmoid, py::arg("x"));
  m.def("tanh",    &ops::tanh,    py::arg("x"));
  m.def("exp",     &ops::exp,     py::arg("x"));
  m.def("log",     &ops::log,     py::arg("x"));
  // ... 60 more
}
\end{lstlisting}

Each call has the form \code{m.def("name", \&cpp\_function,
py::arg("arg") = default)}. The \code{py::arg} machinery preserves
argument names and default values across the binding, which means
the Python-side caller can use keyword arguments and get IDE
autocomplete.

\subsection{Code generation for repetitive bindings}
\label{ssec:python_cpp_binding:codegen}

The unary and binary bindings are repetitive enough that we
generate them from a single declarative table. \code{BindingGen.h}
provides macros that take an op name, an op kind, and a category;
each invocation produces the \code{m.def(...)} line above. The
table is maintained in \code{bind\_ufunc.cpp} itself, near the
top, so the file is self-contained and the generated output
follows immediately.

The alternative (code generation in CMake or in a separate Python
script) was rejected as adding build-system complexity for marginal
benefit.

\section{Exception Translation}
\label{sec:python_cpp_binding:exception_translation}

\Lucid's engine raises a single C++ exception type,
\code{LucidError}, defined in \code{core/Error.h}. The binding
layer translates this exception into a Python exception of the same
name. The translator is installed in \code{bind\_errors.cpp}:

\begin{lstlisting}[style=lucidcpp,language=C++]
void bind_errors(py::module_& m) {
  static py::exception<LucidError> py_lucid_error(m, "LucidError");
  py::register_exception_translator([](std::exception_ptr p) {
    try {
      if (p) std::rethrow_exception(p);
    } catch (const LucidError& e) {
      py_lucid_error(e.what());
    }
  });
}
\end{lstlisting}

The translator runs whenever a C++ function called from Python
throws a \code{LucidError}. The Python side sees a
\code{lucid.LucidError} with the original message, with no further
adaptation needed at the call site.

\subsection{Non-Lucid C++ exceptions}
\label{ssec:python_cpp_binding:other_exceptions}

C++ exceptions other than \code{LucidError} are translated by
pybind11's default mechanism: \code{std::invalid\_argument}
becomes \code{ValueError}, \code{std::out\_of\_range} becomes
\code{IndexError}, \code{std::runtime\_error} becomes
\code{RuntimeError}, and so on. \Lucid's engine code raises
\code{LucidError} for everything that should reach the user; the
default translations are a safety net for programmer errors
(invariant violations, internal asserts) that should not happen in
production.

\subsection{The \texttt{LucidError} message format}
\label{ssec:python_cpp_binding:error_format}

\code{LucidError}'s message is constructed through a fluent builder
(\code{ErrorBuilder} in \code{core/ErrorBuilder.cpp}) that supports
attaching context: the operation name, the tensor shapes and
dtypes, the device, the file and line of the throwing call. The
builder produces structured strings that production users can grep
and that automated diagnostics can parse. The format is documented
in \chref{ch:error_handling}.

\section{The \_C\_\{name\} Alias Convention (H2)}
\label{sec:python_cpp_binding:h2_alias}

Hard Rule H2 (\appref{app:hard_rules}) prescribes a uniform
aliasing convention for every import from \code{lucid.\_C}:

\begin{lucidrule}
Imports from \code{lucid.\_C} must be aliased as
\code{\_C\_\{name\}}, e.g.\ \code{from lucid.\_C import engine as
\_C\_engine}. Bare imports (\code{from lucid.\_C import engine}) and
short aliases (\code{as e}, \code{as eng}, \code{as engine}) are
forbidden.
\end{lucidrule}

The motivation:

\begin{enumerate}
  \item \textbf{Grep-ability.} Every reference to the engine
    binding is prefixed with \code{\_C\_}, so a grep for
    \code{\_C\_} finds every engine call in the codebase. This is
    the single most-used grep when investigating a binding-layer
    bug.
  \item \textbf{Disambiguation.} The \code{\_C\_} prefix signals
    to the reader that the symbol crosses the Python/C++ boundary.
    A bare \code{engine.something()} could be any module; an
    \code{\_C\_engine.something()} cannot be mistaken.
  \item \textbf{Refactor safety.} A future restructuring of the
    binding layer (for example, splitting the single engine
    module into multiple sub-modules) can be applied uniformly
    across the codebase by a search-and-replace on \code{\_C\_engine}.
    Bare aliases would have made the refactor non-mechanical.
\end{enumerate}

The rule is enforced by a static check
(\code{tools/check\_c\_import\_naming.py}) that runs in CI.

\section{GIL Handling}
\label{sec:python_cpp_binding:gil}

Python's Global Interpreter Lock (GIL) is held by default during
every C++ function call from Python. For functions that perform
substantial work without touching Python objects, releasing the
GIL allows other Python threads (e.g., data-loader workers) to run
concurrently.

\Lucid's binding releases the GIL inside operations that:

\begin{itemize}
  \item Dispatch a kernel through \MLX\ or \Accelerate. Kernel
    launches are non-trivial-duration operations that do not touch
    Python.
  \item Perform a \code{backward()} traversal. The autograd engine
    walks the graph without Python interaction.
  \item Block on a CPU-to-GPU bridge or vice versa. The bridge
    waits on \Metal\ events.
\end{itemize}

The GIL release uses pybind11's \code{py::gil\_scoped\_release}:

\begin{lstlisting}[style=lucidcpp,language=C++]
TensorImpl ops::matmul(const TensorImpl& a, const TensorImpl& b) {
  py::gil_scoped_release release;
  return Dispatcher::instance().matmul(a, b);
}
\end{lstlisting}

The scope guard releases the GIL on construction and reacquires on
destruction. The reacquisition is automatic and exception-safe.

\subsection{When not to release}
\label{ssec:python_cpp_binding:no_release}

The GIL is \emph{not} released in functions that:

\begin{itemize}
  \item Touch Python objects (callbacks, exceptions to be raised
    on the calling thread).
  \item Are short enough that the release/reacquire overhead would
    dominate. The break-even point on Apple silicon is roughly
    1\,$\mu$s of useful work; functions shorter than that should
    keep the GIL.
  \item Are called recursively from a context that has already
    released. The scope guard supports nesting but the explicit
    release adds noise without benefit.
\end{itemize}

The rule of thumb: release for any function that dispatches to a
backend kernel; do not release for property getters, type queries,
or other constant-time accessors.

\section{Memory and Reference Counting}
\label{sec:python_cpp_binding:memory}

The interaction between Python's reference counting and C++'s
\code{shared\_ptr} is the subtlest part of the binding. pybind11
handles most of it transparently, but the discipline is worth
stating.

\subsection{shared\_ptr held by Python}
\label{ssec:python_cpp_binding:py_shared_ptr}

A Python \code{Tensor} object holds a \code{shared\_ptr<TensorImpl>}
in its C++ representation. The \code{shared\_ptr}'s reference
count is incremented when the Python object is created and
decremented when the Python object is garbage-collected.

Multiple Python aliases of the same tensor (for example,
\code{a = b}) share the same \code{TensorImpl} through the same
\code{shared\_ptr}; pybind11 handles this by recognising that the
incoming object is already wrapped and returning the existing
Python object rather than constructing a new one.

\subsection{shared\_ptr passed between C++ functions}
\label{ssec:python_cpp_binding:cpp_shared_ptr}

Inside the engine, operations pass \code{shared\_ptr<TensorImpl>}
around freely. The autograd graph holds \code{shared\_ptr}
references to tensors saved for backward; the dispatcher constructs
\code{shared\_ptr}-wrapping return values. The Python side does not
need to know about these intermediate references; the lifetime
they pin is observable only as ``memory usage that goes down after
\code{backward()}.''

\subsection{The risk: cycles}
\label{ssec:python_cpp_binding:cycles}

\code{shared\_ptr} cycles cause leaks: two objects holding
\code{shared\_ptr} to each other never have their reference count
reach zero. The autograd graph could in principle form cycles
(through hooks or through closures), and we are careful in two
specific places: \code{nn.Module}'s parameter registry holds
\code{weak\_ptr}-equivalent references to its children to avoid
parent-child cycles, and the \code{grad\_fn} chain is
unidirectional so that no two backward nodes hold each other.

The discipline is checked at code review; an explicit check would
require a runtime cycle detector that we have not built.

\section{Build System Integration}
\label{sec:python_cpp_binding:build}

The binding layer is built as a CPython extension module using
scikit-build-core, with CMake as the underlying build system. The
configuration files:

\begin{itemize}
  \item \code{pyproject.toml}: declares scikit-build-core as the
    build backend and pins build-time dependencies (pybind11,
    cmake).
  \item \code{CMakeLists.txt}: top-level CMake; declares the
    \code{engine} target and links it against pybind11, \MLX, and
    \Accelerate.
  \item \code{lucid/\_C/CMakeLists.txt}: declares the per-layer
    sub-libraries described in \chref{ch:layered_dag} and links
    them into the engine extension.
\end{itemize}

The build produces a single \code{.so} file (or, on macOS, a
\code{.cpython-XYZ-darwin.so}) that pybind11's loading machinery
recognises and exposes as the \code{engine} module under
\code{lucid.\_C}.

\subsection{The wheel and its constraints}
\label{ssec:python_cpp_binding:wheel}

The published wheel is built once per Python version against the
minimum-target macOS SDK
(\chref{ch:mlx_runtime}~Section~\ref{ssec:mlx_runtime:macos_constraint}).
The wheel tag is \code{cp314-cp314-macosx\_26\_0\_arm64}, which
encodes the Python version (3.14), the platform
(\code{macosx\_26\_0}), and the architecture (\code{arm64}). The
constraints come from \MLX: \Lucid's wheel must match \MLX's wheel
tags to load successfully.

This is the reason \Lucid\ does not currently distribute Linux or
x86 wheels: \MLX\ does not, and \Lucid\ cannot run without \MLX.

\section{Summary and Forward References}
\label{sec:python_cpp_binding:summary}

The Python/C++ binding in \Lucid\ is built on pybind11, organised
into a single \code{engine} module with per-subsystem binding
files, and disciplined by Hard Rule H2 (the \code{\_C\_\{name\}}
alias convention). Exception translation is centralised in
\code{bind\_errors.cpp}, GIL release is per-function with a
1\,$\mu$s rule of thumb, and reference counting interoperates
through \code{shared\_ptr<TensorImpl>} on both sides. The build
produces a single extension module per Python version, tagged to
match \MLX's wheel constraints.

The next chapter, \chref{ch:lazy_imports}, closes
\partref{part:system} with the lazy-import surface that the
top-level Python package exposes. \chref{ch:error_handling}
returns to the exception system from the diagnostics angle.

\begin{lucidnote}
For the user: the binding layer is invisible. The user imports
\code{lucid}, constructs tensors and modules, and never directly
touches the engine module. The only artefact that surfaces is the
\code{LucidError} exception class, which is the single error type
the framework raises.
\end{lucidnote}
